% appD —— 附录 D 微分方程的渐近（原文 wcf-paper.tex 3823–3899 行，PDF 61–62 页；共 4 段、14 条编号公式）
% 原文宏 \plog 不在本工程 preamble 中，此处本地补齐（定义与原文 wcf-paper.tex 159 行相同；ch04 已 providecommand，重复无害）。
% 原文把 \sf 重定义为 \mathrm{sf}（wcf-paper.tex 157 行），本工程对应宏为 \sfm，故原文的 \sf 一律写作 \sfm。
\providecommand{\plog}[1]{#1\partial_{#1}}
\section{微分方程的渐近}
\label{app:diff-asymptotics}

在本附录中，我们希望弄清楚如何计算 $\CX$ 所满足的微分方程的系数。考虑 $\CX$ 在 $\zeta=0$ 附近的渐近展开：
\begin{equation} \label{eq:as-X}
\CX(\theta,\zeta)_\gamma \sim
\CX^{\sfm}_\gamma(\Upsilon_0(\theta),\zeta) \exp \sum_{n>0} \zeta^n
\gamma \cdot g_n(\theta).
\end{equation}
我们考虑第 5.5 节所定义的微分算子，% TODO 合稿改 \ref{sec:differential-problem}（第 5.5 节即原文 §5.5 ``Differential equations''）
\begin{equation}
\frac{1}{\CX_\gamma}\frac{\partial \CX_\gamma}{\partial t^n}  =
\frac{1}{\CX_\gamma} \CA_n \CX_\gamma.
\end{equation}
把展开式 \eqref{eq:as-X} 代入，只保留最初几项：
\begin{equation}
\frac{\partial \log \CX^{\sfm}_\gamma(\Upsilon_0(\theta))}{\partial
t^n} = \left(\frac{1}{\zeta}\CA^{(-1)}_n + \CA^{(0)}_n\right)
\Upsilon_0(\theta)_\gamma+ \CA_n^{(-1)} \gamma \cdot g_1(\theta).
\end{equation}
$\zeta$ 的领头部分是一条我们已经理解的陈述：
\begin{equation}
\pi R \frac{\partial Z_\gamma}{\partial t^n} = \CA^{(-1)}_n
\Upsilon_0(\theta)_\gamma.
\end{equation}
这意味着，$\CA^{(-1)}$ 是 $\CA^{(-1)}_{\sfm}$ 经 $\Upsilon_0$ 拉回所得。

展开中的下一项是
\begin{equation}
i \frac{\partial \Upsilon_0(\theta)_\gamma}{\partial t^n} =
\CA^{(0)}_n \Upsilon_0(\theta)_\gamma+ \CA_n^{(-1)} \gamma \cdot
g_1(\theta),
\end{equation}
在已知 $\Upsilon_0(\theta)$ 的前提下，它确定了 $\CA^{(0)}_n$。

还有另一种观点，它相当有用：考虑各个微分方程之间的相容性条件。例如，考虑方程 $[R\partial_R - \CA_R, \zeta \partial_\zeta - \CA_\zeta] = 0$，亦即
\begin{equation} \left[ \plog{R} + \frac{1}{\zeta} {\CA_{\zeta}^{(-1)}} - \CA_R^{(0)} -
\zeta \CA^{(1)}_{\zeta} \, , \, \plog{\zeta} -
\frac{1}{\zeta} {\CA_{\zeta}^{(-1)}} - \CA_{\zeta}^{(0)} - \zeta
\CA^{(1)}_{\zeta} \right] = 0,
\end{equation}
把它按 $\zeta$ 的幂次展开。这给出三个方程：
\begin{align}
\plog{R} \CA_{\zeta}^{(-1)} - [\CA_R^{(0)},\CA_{\zeta}^{(-1)}] &=
[\CA_{\zeta}^{(0)},\CA_{\zeta}^{(-1)}] + \CA_{\zeta}^{(-1)},\\
\plog{R} \CA_{\zeta}^{(0)} - [\CA_R^{(0)},\CA_{\zeta}^{(0)}] &=
2 [\CA_{\zeta}^{(1)},\CA_{\zeta}^{(-1)}],\\
\plog{R} \CA_{\zeta}^{(1)} - [\CA_R^{(0)},\CA_{\zeta}^{(1)}] &=
[\CA_{\zeta}^{(1)},\CA_{\zeta}^{(0)}] + \CA_{\zeta}^{(1)}.
\end{align}
这些方程与 Nahm 方程极为相像，与后者的差别仅在于右端多出的两个线性项。这两个额外项在大半径处占主导。要推导大 $R$ 渐近，另一种策略仍是围绕半平坦解对这三个方程作迭代求解。

另一组有趣的「等单值」方程，可以通过类似地展开 $[\partial_u - \CA_u, \zeta\partial_\zeta - \CA_\zeta] = 0$ 得到：
\begin{align}
0 &=[\CA^{(-1)}_u, \CA^{(-1)}_\zeta], \\
\frac{\partial}{\partial u} \CA_{\zeta}^{(-1)} -
[\CA_u^{(0)},\CA_{\zeta}^{(-1)}] &=
[\CA_u^{(-1)},\CA_{\zeta}^{(0)}] - \CA_{u}^{(-1)},\\
\frac{\partial}{\partial u} \CA_{\zeta}^{(0)} -
[\CA_u^{(0)},\CA_{\zeta}^{(0)}]  &=
[\CA_u^{(-1)},\CA_{\zeta}^{(1)}], \\
\frac{\partial}{\partial u} \CA_{\zeta}^{(1)} -
[\CA_u^{(0)},\CA_{\zeta}^{(1)}] &=
[\CA_u^{(1)},\CA_{\zeta}^{(0)}] + \CA_{u}^{(-1)}, \\
0 &= [\CA^{(1)}_u, \CA^{(1)}_\zeta].
\end{align}
